% Part: lambda-calculus
% Chapter: syntax
% Section: eta

\documentclass[../../../include/open-logic-section]{subfiles}

\begin{document}

\olfileid{lam}{syn}{eta}
\olsection{$\eta$-روپ بدلنا}

$\lambda$ ٹرماں اُتے اک ہور تعلق اے۔ \olref[fv]{sec} وچ
اسیں $\lambd[x][(fx)]$ دی مثال ورتی سی، جیہڑا اک دلیل قبول
کردا اے تے اوہدے اُتے $f$ لاؤندا اے۔ دوجے لفظاں وچ ایہ اوہی
فنکشن اے جو~$f$ اے۔ $\lambd[x][(fx)]N$ تے $fN$ دوہاں دی
تقلیل نال~$fN$ ملدا اے۔ اسیں ایس خیال نوں $\eta$-تقلیل
تے $\eta$-وادھے نال ظاہر کردے آں۔

\begin{defn}[$\eta$-سکڑاؤ، $\eredone$]
  \ollabel{defn:beredone}
  \emph{$\eta$-سکڑاؤ} ($\eredone$) ٹرماں اُتے سبھ توں چھوٹا
  موافق تعلق اے جیہڑا ہِٹھلی شرط پوری کردا اے۔
  \[
    \lambd[x][M x] \eredone M \text{ بشرطیکہ } x \notin FV(M)
  \]
\end{defn}

\begin{defn}[$\beta\eta$-تقلیل، $\bered$] \ollabel{defn:bered}
  \emph{$\beta\eta$-تقلیل} ($\bered$) ٹرماں اُتے سبھ توں چھوٹا
  انعکاسی تے تعدّی والا تعلق اے جیہدے وچ $\bredone$ تے
  $\eredone$ شامل نیں۔ یعنی ایہدے قاعدے انعکاس تے تعدّی دے
  قاعدے تے ایہ دو ہور قاعدے نیں۔
  \begin{enumerate}
  \item جے $M \bredone N$، تاں $M \bered N$۔ \ollabel{defn:bered3}
  \item جے $M \eredone N$، تاں $M \bered N$۔ \ollabel{defn:bered4}
  \end{enumerate}

\end{defn}

\begin{defn}
  اسیں ہم ارزیتی تعلق $\equal$ نوں $\eta$-روپ بدلن دے قاعدے نال
  ودھاؤندے آں۔
  \[
  \lambd[x][f x] \equal f
  \]
  ودھے ہوئے تعلق نوں $\equal[\eta]$ لکھدے آں۔
\end{defn}

$\eta$-ہم ارزی اہم اے، کیوں جے ایہ لامبڈا ٹرماں دی
امتدادی برابری نال جُڑی اے۔

\begin{defn}[امتدادی برابری]
  اسیں ہم ارزیتی تعلق $\equal$ نوں (\ext) دے قاعدے نال ودھاؤندے آں۔
  \begin{center}
  جے $Mx \equal Nx$، تاں $M \equal N$، بشرطیکہ $x \notin FV(MN)$۔
  \end{center}
  ودھے ہوئے تعلق نوں $\equal[\ext]$ لکھدے آں۔
\end{defn}

موٹے طور اُتے قاعدہ ایہ آکھدا اے کہ جے دو ٹرم فنکشن
وجوں اکّو دلیل اُتے اکّو جیہا ورتاؤ کرن، تاں اوہناں نوں
برابر سمجھنا چاہیدا اے۔

ہُن ثابت کردے آں کہ $\eta$ دا قاعدہ عین امتدادی برابری
دیندا اے تے ایس توں ودھ کجھ نہیں دیندا۔

\begin{thm}
  $M \equal[\ext] N$ اُدوں تے صرف اُدوں اے جدوں $M \equal[\eta] N$۔
\end{thm}

\begin{proof}
  پہلاں اسیں ثابت کردے آں کہ $\equal[\eta]$ امتدادی برابری
  دے قاعدے ہِٹھ بند اے۔ یعنی $ext$ دا قاعدہ $\equal[\eta]$ وچ
  کجھ ہور نہیں پاؤندا۔ ایس لئی $\equal[\eta]$ وچ $\equal[\ext]$
  شامل اے، تے جے $M \equal[\ext] N$، تاں $M \equal[\eta] N$۔

  ایہ ثابت کرن لئی کہ $\equal[\eta]$، \ext ہِٹھ بند اے،
  دھیان دیو کہ جے کوئی $M
  \equal N$، \ext{} دے قاعدے نال اخذ ہویا اے، تاں اوہدا
  مفروضہ $Mx
  \equal[\eta] Nx$ اے۔ فیر $\lambd[x][Mx]
  \equal[\eta] \lambd[x][Nx]$ ملدا اے $\equal$ دے اک قاعدے
  نال۔ دوہاں پاسیاں اُتے $\eta$ لاؤن نال $M \equal[\eta] N$
  ملدا اے۔

  ایس طرح ثابت کردے آں کہ $\eta$ دا قاعدہ $\equal[ext]$ وچ
  شامل اے۔ ہر $\lambd[x][Mx]$ تے $M$ لئی، جتھے $x \notin
  FV(M)$، ساڈے کول $(\lambd[x][Mx])x \equal[\ext] Mx$ اے۔
  \ext{} دے قاعدے نال ایس توں $\lambd[x][Mx] \equal[\ext] M$
  ملدا اے۔
\end{proof}

\end{document}
